\documentclass[12pt,a4paper]{article}
\usepackage[utf8]{vietnam}
\usepackage[top=2cm, bottom=2cm, left=1.5cm, right=1cm]{geometry}
\usepackage{amsmath,amsfonts,amssymb}
\usepackage{tabularx,longtable,multirow}
\usepackage{array}
\usepackage{fancyhdr}
\usepackage{enumitem}
\usepackage{xcolor}

% Cấu hình cột bảng để không bị tràn lề
\newcolumntype{L}[1]{>{\raggedright\arraybackslash}p{#1}}
\newcolumntype{C}[1]{>{\centering\arraybackslash}p{#1}}

% Cấu hình Header/Footer
\pagestyle{fancy}
\fancyhf{}
\rhead{\small GV: Nguyễn Văn Sang}
\lhead{\small Kế hoạch giáo dục 2025-2026 - Môn Toán 12}
\cfoot{\thepage}

\begin{document}

% --- PHẦN TIÊU NGỮ ---
\noindent
\begin{minipage}[t]{0.5\textwidth}
    \textbf{TRƯỜNG: THPT NGUYỄN HỮU CẢNH}\\
    \textbf{TỔ: Toán}\\
    \textbf{Họ và tên giáo viên: Nguyễn Văn Sang}
\end{minipage}
\begin{minipage}[t]{0.5\textwidth}
    \centering
    \textbf{CỘNG HÒA XÃ HỘI CHỦ NGHĨA VIỆT NAM}\\
    \textbf{Độc lập - Tự do - Hạnh phúc}\\
    ----------------------------
\end{minipage}

\vspace{0.8cm}
\begin{center}
    \Large \textbf{KẾ HOẠCH GIÁO DỤC CỦA GIÁO VIÊN} \\
    \large \textbf{(Năm học 2025 - 2026)}
\end{center}

\section*{I. THÔNG TIN CÁ NHÂN}
\begin{itemize}[noitemsep]
    \item Họ và tên: Nguyễn Văn Sang \quad | \quad Năm sinh: 1987
    \item Năm vào ngành: 2010 \quad | \quad Chức vụ: Giáo viên
    \item Giảng dạy: Toán \quad | \quad Công tác kiêm nhiệm: Không
\end{itemize}

\section*{II. KẾ HOẠCH CÁ NHÂN}
\subsection*{1. Những căn cứ xây dựng kế hoạch}
\begin{itemize}[label=-, noitemsep]
    \item Công văn 4612/BGDĐT-GDTrH ngày 03/10/2017 về hướng dẫn chương trình GDPT.
    \item Quyết định số 3089/QĐ-UBND ngày 08/8/2025 của UBND TP. Hồ Chí Minh.
    \item Văn bản số 4977/SGDĐT-GDTrH ngày 13/8/2025 của Sở Giáo dục và Đào tạo TP. HCM.
    \item Kế hoạch số 4818/KH-SGDĐT ngày 07/8/2025 về bồi dưỡng thường xuyên.
    \item Căn cứ vào năng lực cá nhân và yêu cầu nhiệm vụ được giao.
\end{itemize}

\subsection*{2. Mục tiêu chung}
- Hoàn thành tốt nhiệm vụ được giao. Đáp ứng Chuẩn nghề nghiệp, phát triển bản thân.

\subsection*{3. Nội dung}
\textbf{3.1. Đặc điểm tình hình:}\\
- \textit{Thuận lợi:} Có tinh thần trách nhiệm cao, nhiệt tình. Trình độ chuyên môn vững vàng, ứng dụng tốt CNTT.\\
- \textit{Khó khăn:} CSVC còn thiếu thốn. Cả 3 khối 10, 11, 12 đều học chương trình mới nên nhiều thứ cần chuẩn bị.

\textbf{3.1.2 Công việc được giao:}\\
Giảng dạy môn Toán các lớp: \textbf{12C1, 12C9}.

\textbf{3.2.2 Giáo dục đạo đức, tư tưởng:}\\
Giáo dục HS lòng yêu nước, tôn trọng pháp luật, nhân ái và tự trọng.

\newpage
% --- III. PHÂN PHỐI CHƯƠNG TRÌNH CHI TIẾT ---
\section*{III. KẾ HOẠCH DẠY HỌC (PHÂN PHỐI CHƯƠNG TRÌNH)}
\textbf{HỌC KỲ I}
\begin{longtable}{|C{1.3cm}|C{0.5cm}|L{4.2cm}|C{0.5cm}|L{4.2cm}|C{0.5cm}|L{2.3cm}|}
\hline
\textbf{Tuần} & \textbf{T} & \textbf{Giải tích} & \textbf{T} & \textbf{HH và Đo lường} & \textbf{T} & \textbf{Chuyên đề} \\ \hline
\endhead

% TUẦN 1
\multirow{2}{*}{\begin{tabular}[c]{@{}c@{}}1\\(8/9-\\14/9)\end{tabular}} & 1 & \textbf{Tính đơn điệu và cực trị của hàm số.} \newline \textit{Yêu cầu:} Nhận biết tính đơn điệu dựa vào đạo hàm; thể hiện qua BBT; nhận biết điểm cực trị; ứng dụng tối ưu hóa lợi nhuận, quỹ đạo bay. & 1 & \textbf{Vectơ và các phép toán trong không gian.} \newline \textit{Yêu cầu:} Nhận biết vectơ, các phép toán (tổng, hiệu, tích vô hướng). Tính và chứng minh đẳng thức. & 1 & \textbf{Bài toán quy hoạch tuyến tính.} \newline \textit{Yêu cầu:} Giải bài toán quy hoạch tuyến tính bằng hệ BPT bậc nhất. \\ \cline{2-3}
 & 2 & Tính đơn điệu và cực trị (tiếp). & & & & \\ \hline

% TUẦN 2
\multirow{2}{*}{\begin{tabular}[c]{@{}c@{}}2\\(15/9-\\21/9)\end{tabular}} & 3 & \textbf{Tính đơn điệu và cực trị (tiếp).} \newline \textit{Yêu cầu:} Vận dụng giải các bài toán liên quan. & 2 & Vectơ và các phép toán (tiếp). & 2 & Bài toán quy hoạch tuyến tính (tiếp). \\ \cline{2-3}
 & 4 & Tính đơn điệu và cực trị (tiếp). & & & & \\ \hline

% TUẦN 3
\multirow{2}{*}{\begin{tabular}[c]{@{}c@{}}3\\(22/9-\\28/9)\end{tabular}} & 5 & \textbf{Tính đơn điệu và cực trị (tiếp).} \newline \textbf{GTLN \& GTNN của hàm số.} \newline \textit{Yêu cầu:} Xác định GTLN, GTNN trên tập cho trước bằng đạo hàm. Ứng dụng thực tiễn. & 3 & Vectơ và các phép toán (tiếp). & 3 & Bài toán quy hoạch tuyến tính (tiếp). \\ \cline{2-7}
 & 6 & GTLN \& GTNN (tiếp). & & & & \\ \hline

% TUẦN 4
\multirow{2}{*}{\begin{tabular}[c]{@{}c@{}}4\\(29/9-\\5/10)\end{tabular}} & 7 & \textbf{GTLN \& GTNN của hàm số. \color{red}KTTX1} \newline \textit{Yêu cầu:} Giải bài tập vận dụng. & 4 & Vectơ và các phép toán (tiếp). & 4 & \textbf{Vận dụng đạo hàm giải bài toán tối ưu.} \newline \textit{Yêu cầu:} Giải bài toán thực tế (khoảng cách, thời gian, kinh tế). \\ \cline{2-3}
 & 8 & GTLN \& GTNN (tiếp). & & & & \\ \hline

% TUẦN 5
\multirow{2}{*}{\begin{tabular}[c]{@{}c@{}}5\\(6/10-\\12/10)\end{tabular}} & 9 & \textbf{Đường tiệm cận của đồ thị hàm số.} \newline \textit{Yêu cầu:} Nhận biết hình ảnh tiệm cận ngang, đứng, xiên. Biết cách tìm tiệm cận. Giải quyết bài toán thực tế. & 5 & \textbf{Tọa độ của vectơ trong không gian.} \newline \textit{Yêu cầu:} Nhận biết tọa độ vectơ đối với hệ trục tọa độ. Xác định độ dài vectơ; biểu thức tọa độ các phép toán. Vận dụng vào thực tiễn. & 5 & \textbf{Vận dụng đạo hàm giải bài toán tối ưu.} \newline \textit{Yêu cầu:} Giải bài toán thực tế tìm GTLN, GTNN trong kinh tế, kỹ thuật. \\ \cline{2-3}
 & 10 & Đường tiệm cận (tiếp). & & & & \\ \hline

% TUẦN 6
\multirow{2}{*}{\begin{tabular}[c]{@{}c@{}}6\\(13/10-\\19/10)\end{tabular}} & 11 & \textbf{Khảo sát và vẽ đồ thị hàm số cơ bản.} \newline \textit{Yêu cầu:} Mô tả sơ đồ tổng quát khảo sát hàm số. Khảo sát các hàm $y=ax^3+bx^2+cx+d$, $y=\frac{ax+b}{cx+d}$, $y=\frac{ax^2+bx+c}{mx+n}$. Nhận biết tính đối xứng. & 6 & Tọa độ của vectơ trong không gian. \newline \textit{Yêu cầu:} Vận dụng định nghĩa tọa độ vectơ để giải các bài toán hình học cơ bản. & 6 & Vận dụng đạo hàm giải bài toán tối ưu. \newline \textit{Yêu cầu:} Giải bài toán thực tế liên quan đến chi phí, doanh thu. \\ \cline{2-3}
 & 12 & Khảo sát và vẽ đồ thị (tiếp). & & & & \\ \hline

% TUẦN 7
\multirow{2}{*}{\begin{tabular}[c]{@{}c@{}}7\\(20/10-\\26/10)\end{tabular}} & 13 & Khảo sát và vẽ đồ thị (tiếp). \newline \textit{Yêu cầu:} Rèn luyện kỹ năng vẽ đồ thị và nhận dạng đồ thị hàm số. & 7 & Tọa độ của vectơ trong không gian. \newline \textit{Yêu cầu:} Tính khoảng cách, góc trong không gian bằng tọa độ. & 7 & Vận dụng đạo hàm giải bài toán tối ưu. \newline \textit{Yêu cầu:} Tổng hợp các bài toán vận dụng cao về cực trị thực tế. \\ \cline{2-3}
 & 14 & Khảo sát và vẽ đồ thị (tiếp). & & & & \\ \hline

% TUẦN 8
\multirow{2}{*}{\begin{tabular}[c]{@{}c@{}}8\\(27/10-\\2/11)\end{tabular}} & 15 & \textbf{Bài tập cuối chương I.} & 8 & Tọa độ của vectơ trong không gian. \newline \textbf{KTTX2} & 8 & \textbf{Bài tập cuối chuyên đề 1.} \\ \cline{2-3}
 & 16 & Bài tập cuối chương I (tiếp). & & & & \\ \hline

\multicolumn{7}{|c|}{\textbf{Kiểm tra giữa học kỳ 1 (TẠI LỚP)}} \\
\multicolumn{7}{|c|}{\textbf{Điểm KT GHK1: HS2}. Hình thức: 3 dạng (CTM). Thời lượng: 90 phút.} \\ \hline

% TUẦN 9
\multirow{2}{*}{\begin{tabular}[c]{@{}c@{}}9\\(3/11-\\9/11)\end{tabular}} & 17 & \textbf{Khoảng biến thiên và khoảng tứ phân vị mẫu ghép nhóm.} \newline \textit{Yêu cầu:} Nhận biết định nghĩa, cách tính $R, IQR$. Hiểu ý nghĩa độ phân tán. So sánh mức độ phân tán. Giải thích mối quan hệ. & 9 & \textbf{Biểu thức tọa độ của các phép toán vectơ.} \newline \textit{Yêu cầu:} Nhận biết các công thức. Xác định biểu thức tọa độ. Vận dụng giải toán liên quan thực tiễn. & 9 & Bài tập cuối chuyên đề 1. \newline \textit{Yêu cầu:} Củng cố kỹ năng giải bài toán quy hoạch tuyến tính và tối ưu hóa. \\ \cline{2-3}
 & 18 & Khoảng biến thiên và khoảng tứ phân vị (tiếp). & & & & \\ \hline

% TUẦN 10
\multirow{2}{*}{\begin{tabular}[c]{@{}c@{}}10\\(10/11-\\16/11)\end{tabular}} & 19 & Khoảng biến thiên và khoảng tứ phân vị (tiếp). \newline \textbf{Phương sai và độ lệch chuẩn mẫu ghép nhóm.} \newline \textit{Yêu cầu:} Nhận biết định nghĩa, cách tính $S^2, S$. & 10 & \textbf{Biểu thức tọa độ các phép toán vectơ (tiếp).} \newline \textit{Yêu cầu:} Vận dụng biểu thức tọa độ tích vô hướng giải bài toán góc và khoảng cách. & 10 & \textbf{Tiền tệ. Lãi suất.} \newline \textit{Yêu cầu:} Nhận biết một số vấn đề tiền tệ. Thiết lập kế hoạch tài chính cá nhân. \\ \cline{2-3}
 & 20 & Phương sai và độ lệch chuẩn (tiếp). & & & & \\ \hline

% TUẦN 11
\multirow{2}{*}{\begin{tabular}[c]{@{}c@{}}11\\(17/11-\\23/11)\end{tabular}} & 21 & Phương sai và độ lệch chuẩn mẫu ghép nhóm (tiếp). & 11 & Biểu thức tọa độ của các phép toán vectơ. \newline \textit{Yêu cầu:} Vận dụng định lý cosin và biểu thức tọa độ để tính góc giữa hai vectơ. & 11 & Tiền tệ. Lãi suất. \newline \textit{Yêu cầu:} Giải bài toán lãi đơn, lãi kép trong các tình huống gửi tiết kiệm. \\ \cline{2-3}
 & 22 & Phương sai và độ lệch chuẩn mẫu ghép nhóm (tiếp). \newline \textit{Yêu cầu:} Tổng hợp và giải quyết các bài tập về các đặc trưng đo độ phân tán. & & & & \\ \hline

% TUẦN 12
\multirow{2}{*}{\begin{tabular}[c]{@{}c@{}}12\\(24/11-\\30/11)\end{tabular}} & 23 & Bài tập cuối chương III. \textbf{\color{red}KTTX3} \newline \textit{Yêu cầu:} Củng cố kiến thức về thống kê mẫu ghép nhóm qua bài tập thực hành. & 12 & Biểu thức tọa độ các phép toán vectơ (tiếp). & 12 & Tiền tệ. Lãi suất (tiếp). \newline \textit{Yêu cầu:} Tìm hiểu về tỷ lệ lạm phát và ảnh hưởng đến giá trị tiền tệ. \\ \cline{2-3}
 & & & 13 & Bài tập cuối chương II. \newline \textit{Yêu cầu:} Hệ thống hóa kiến thức về vectơ và tọa độ trong không gian. & & \\ \hline

% TUẦN 13
\multirow{2}{*}{\begin{tabular}[c]{@{}c@{}}13\\(1/12-\\7/12)\end{tabular}} & 24 & Ôn tập kiểm tra cuối học kỳ 1. & 14 & Ôn tập kiểm tra cuối học kỳ 1. & 13 & Tiền tệ. Lãi suất (tiếp). \newline \textit{Yêu cầu:} -nt \\ \cline{2-3}
 & 25 & Ôn tập kiểm tra cuối học kỳ 1. \newline \textit{Yêu cầu:} Thống nhất theo ma trận đề. & & \textit{Yêu cầu:} Thống nhất theo ma trận đề. & & \\ \hline

% TUẦN 14
\multirow{2}{*}{\begin{tabular}[c]{@{}c@{}}14\\(8/12-\\14/12)\end{tabular}} & 26 & Ôn tập kiểm tra cuối học kỳ 1. & 15 & Ôn tập kiểm tra cuối học kỳ 1. & 14 & \textbf{Tín dụng. Vay nợ.} \newline \textit{Yêu cầu:} Nhận biết kết quả trả nợ đúng hạn; hồ sơ tín dụng và giá trị tín dụng. Vận dụng tính toán kế hoạch trả nợ. \\ \cline{2-3}
 & 27 & Ôn tập kiểm tra cuối học kỳ 1. \newline \textit{Yêu cầu:} Giải các dạng toán trắc nghiệm tổng hợp theo đề minh họa. & & \textit{Yêu cầu:} Tái hiện các công thức hình học tọa độ quan trọng. & & \\ \hline

\multicolumn{7}{|c|}{\textbf{Kiểm tra cuối học kỳ 1 (TẬP TRUNG)}} \\
\multicolumn{7}{|c|}{Hình thức: 3 dạng (CTM). Thời lượng làm bài: 90 phút.} \\ \hline

% TUẦN 15-16
\multirow{2}{*}{\begin{tabular}[c]{@{}c@{}}15-16 \\ (15/12 - 28/12)\end{tabular}} & & \textbf{Chấm bài và sửa bài kiểm tra cuối học kỳ 1.} & & & & \\ \cline{2-3}
 & & & & & & \\ \hline
\end{longtable}

\textbf{HỌC KỲ II}
\begin{longtable}{|C{1.3cm}|C{0.5cm}|L{4.2cm}|C{0.5cm}|L{4.2cm}|C{0.5cm}|L{2.3cm}|}
\hline
\textbf{Tuần} & \textbf{T} & \textbf{Giải tích} & \textbf{T} & \textbf{HH và Đo lường} & \textbf{T} & \textbf{Chuyên đề} \\ \hline
\endhead

% TUẦN 1
\multirow{2}{*}{\begin{tabular}[c]{@{}c@{}}1\\(19/01-\\25/01)\end{tabular}} & 3 & \textbf{Nguyên hàm (tiếp).} \newline \textit{Yêu cầu:} Nhận biết khái niệm, tính chất cơ bản; xác định nguyên hàm một số hàm sơ cấp. Tính nguyên hàm trong trường hợp đơn giản. & 2 & \textbf{Phương trình mặt phẳng (tiếp).} \newline \textit{Yêu cầu:} Nhận biết PT tổng quát; thiết lập PT mặt phẳng; tính khoảng cách từ điểm đến mặt phẳng. & 17 & \textbf{Đầu tư tài chính. Lập kế hoạch tài chính cá nhân.} \newline \textit{Yêu cầu:} Nhận biết vấn đề đầu tư; vận dụng kiến thức giải quyết vấn đề. \\ \cline{2-3}
 & 4 & Nguyên hàm (tiếp). & & & & \\ \hline

% TUẦN 2
\multirow{2}{*}{\begin{tabular}[c]{@{}c@{}}2\\(26/01-\\01/02)\end{tabular}} & 5 & Nguyên hàm (tiếp). & 3 & Phương trình mặt phẳng (tiếp). \newline \textit{Yêu cầu:} Vận dụng điều kiện song song, vuông góc của hai mặt phẳng. & 18 & Đầu tư tài chính (tiếp). \newline \textit{Yêu cầu:} Phân tích rủi ro và lợi nhuận của các danh mục đầu tư. \\ \cline{2-3}
 & 6 & \textbf{Tích phân.} \newline \textit{Yêu cầu:} Nhận biết định nghĩa và tính chất của tích phân. Giải quyết bài toán liên quan thực tiễn. & & & & \\ \hline

% TUẦN 3
\multirow{2}{*}{\begin{tabular}[c]{@{}c@{}}3\\(02/02-\\08/02)\end{tabular}} & 7 & \textbf{Tích phân (tiếp).} \newline \textit{Yêu cầu:} Sử dụng các phương pháp tính tích phân cơ bản (đổi biến, từng phần đơn giản). & 4 & \textbf{Phương trình đường thẳng.} \newline \textit{Yêu cầu:} Nhận biết phương trình chính tắc, tham số, vectơ chỉ phương của đường thẳng trong không gian. & 19 & Đầu tư tài chính (tiếp). \newline \textit{Yêu cầu:} Xây dựng mô hình bài toán đầu tư đơn giản. \\ \cline{2-3}
 & 8 & Tích phân (tiếp). & & & & \\ \hline

% TUẦN 4
\multirow{2}{*}{\begin{tabular}[c]{@{}c@{}}4\\(09/02-\\15/02)\end{tabular}} & 9 & Tích phân (tiếp). \newline \textbf{KTTX1} \newline \textit{Yêu cầu:} Kiểm tra khả năng vận dụng tính chất tích phân giải bài tập. & 5 & Phương trình đường thẳng (tiếp). \newline \textit{Yêu cầu:} Thiết lập PT đường thẳng; xác định góc giữa hai đường thẳng, đường thẳng và mặt phẳng. & 20 & \textbf{Bài tập cuối chuyên đề 2.} \newline \textit{Yêu cầu:} Hệ thống hóa kiến thức về toán học trong tài chính. \\ \cline{2-3}
 & 10 & Tích phân (tiếp). & & & & \\ \hline

% TUẦN 5
\multirow{2}{*}{\begin{tabular}[c]{@{}c@{}}5\\(23/02-\\01/03)\end{tabular}} & 11 & \textbf{Ứng dụng hình học của tích phân.} \newline \textit{Yêu cầu:} Tính diện tích hình phẳng, thể tích khối tròn xoay. Vận dụng giải bài toán thực tiễn. & 6 & Phương trình đường thẳng (tiếp). \newline \textit{Yêu cầu:} Tính khoảng cách giữa hai đường thẳng chéo nhau. & 21 & Bài tập cuối chuyên đề 2 (tiếp). \newline \textit{Yêu cầu:} Giải bài toán ứng dụng chuyên đề tài chính. \\ \cline{2-3}
 & 12 & Ứng dụng hình học của tích phân (tiếp). & & & & \\ \hline

% TUẦN 6
\multirow{2}{*}{\begin{tabular}[c]{@{}c@{}}6\\(02/03-\\08/03)\end{tabular}} & 13 & Ứng dụng hình học của tích phân (tiếp). \newline \textit{Yêu cầu:} Giải quyết bài tập về thể tích các vật thể trong thực tế. & 7 & \textbf{Phương trình mặt cầu. KTTX2} \newline \textit{Yêu cầu:} Nhận biết PT mặt cầu; xác định tâm, bán kính; thiết lập PT mặt cầu khi biết tâm và bán kính. & 22 & \textbf{Biến ngẫu nhiên rời rạc.} \newline \textit{Yêu cầu:} Nhận biết khái niệm; tính kì vọng, phương sai, độ lệch chuẩn. \\ \cline{2-3}
 & 14 & Ứng dụng hình học của tích phân (tiếp). & & & & \\ \hline

% TUẦN 7
\multirow{2}{*}{\begin{tabular}[c]{@{}c@{}}7\\(09/03-\\15/03)\end{tabular}} & 15 & \textbf{Bài tập cuối chương IV.} \newline \textit{Yêu cầu:} Củng cố kiến thức về nguyên hàm, tích phân và ứng dụng. & 8 & Phương trình mặt cầu (tiếp). \newline \textit{Yêu cầu:} Vận dụng phương trình mặt cầu vào các bài toán cực trị. & 23 & Biến ngẫu nhiên rời rạc (tiếp). \newline \textit{Yêu cầu:} Ý nghĩa và ứng dụng của kỳ vọng trong thực tế. \\ \cline{2-3}
 & 16 & & 9 & & & \\ \hline

\multicolumn{7}{|c|}{\textbf{Kiểm tra giữa học kỳ 2 (TẠI LỚP)}} \\
\multicolumn{7}{|c|}{Hình thức: 3 dạng (CTM). Thời lượng làm bài: 90 phút.} \\ \hline

% TUẦN 8
\multirow{2}{*}{\begin{tabular}[c]{@{}c@{}}8\\(16/03-\\22/03)\end{tabular}} & 17 & \textbf{Xác suất có điều kiện.} \newline \textit{Yêu cầu:} Nhận biết khái niệm; ý nghĩa của xác suất có điều kiện trong tình huống thực tiễn. & 10 & \textbf{Bài tập cuối chương V.} \newline \textit{Yêu cầu:} Hệ thống hóa kiến thức hình học tọa độ trong không gian. & 24 & \textbf{KTTX3} \newline \textbf{Phép thử lặp và công thức Bernoulli.} \newline \textit{Yêu cầu:} Nhận biết khái niệm; vận dụng giải bài toán thực tiễn. \\ \cline{2-3}
 & 18 & Xác suất có điều kiện (tiếp). & & & & \\ \hline

% TUẦN 9
\multirow{2}{*}{\begin{tabular}[c]{@{}c@{}}9\\(23/03-\\29/03)\end{tabular}} & 19 & \textbf{Công thức xác suất toàn phần và công thức Bayes.} \newline \textit{Yêu cầu:} Mô tả công thức thông qua bảng dữ liệu thống kê 2x2 và sơ đồ hình cây. & 11 & \textbf{Luyện tập hình học Oxyz.} \newline \textit{Yêu cầu:} Tổng ôn phương trình mặt phẳng, đường thẳng, mặt cầu. & 25 & \textbf{Phân bố Bernoulli và phân bố nhị thức.} \newline \textit{Yêu cầu:} Nhận biết khái niệm; vận dụng giải toán. \\ \cline{2-3}
 & 20 & Công thức xác suất toàn phần và Bayes (tiếp). & & & & \\ \hline

% TUẦN 10
\multirow{2}{*}{\begin{tabular}[c]{@{}c@{}}10\\(30/03-\\05/04)\end{tabular}} & 21 & \textbf{Bài tập cuối chương VI.} \newline \textbf{Tính giá trị gần đúng tích phân bằng máy tính cầm tay.} \newline \textit{Yêu cầu:} Sử dụng máy tính hỗ trợ kiểm tra kết quả. & 11 & \textbf{Sử dụng phần mềm GeoGebra để biểu diễn hình học tọa độ trong không gian.} \newline \textit{Yêu cầu:} Nhận biết và vẽ các đối tượng Oxyz. & 26 & \textbf{Phân bố Bernoulli và phân bố nhị thức (tiếp).} \newline \textit{Yêu cầu:} Vận dụng tính xác suất trong các bài toán thực tế phức tạp. \\ \cline{2-3}
 & 22 & (Tiếp theo) \newline \textbf{KTTX3} & & & & \\ \hline

% TUẦN 11
\multirow{2}{*}{\begin{tabular}[c]{@{}c@{}}11\\(06/04-\\12/04)\end{tabular}} & 23 & \textbf{Minh họa và tính tích phân bằng phần mềm GeoGebra.} \newline \textit{Yêu cầu:} Trực quan hóa phần hình phẳng và vật thể tròn xoay. & 12 & \textbf{Ôn thi cuối học kỳ 2.} \newline \textit{Yêu cầu:} Thống nhất theo ma trận đề. & 26 & Phân bố Bernoulli và nhị thức (tiếp). \newline \textit{Yêu cầu:} Củng cố kiến thức xác suất chương IV. \\ \cline{2-3}
 & 24 & \textbf{Ôn thi cuối học kỳ 2.} \newline \textit{Yêu cầu:} Luyện đề trắc nghiệm tổng hợp. & & & & \\ \hline

% TUẦN 12
\multirow{2}{*}{\begin{tabular}[c]{@{}c@{}}12\\(13/04-\\19/04)\end{tabular}} & 25 & \textbf{Ôn thi cuối học kỳ 2.} & 13 & \textbf{Ôn thi cuối học kỳ 2 (tiếp).} \newline \textit{Yêu cầu:} Giải quyết các câu hỏi vận dụng và vận dụng cao. & 27 & \textbf{Bài tập cuối chuyên đề 3.} \newline \textit{Yêu cầu:} Hệ thống hóa kiến thức xác suất nhị thức và ứng dụng. \\ \cline{2-3}
 & 26 & \textbf{Ôn thi cuối học kỳ 2 (tiếp).} \newline \textit{Yêu cầu:} Hoàn thiện kỹ năng trình bày và làm bài. & & & & \\ \hline

% TUẦN 13 - 14
\multicolumn{7}{|c|}{\textbf{KIỂM TRA CUỐI HỌC KỲ 2 (TẬP TRUNG)}} \\
\multicolumn{7}{|c|}{Hình thức: 3 dạng (CTM). Thời lượng làm bài: 90 phút.} \\ \hline
\multicolumn{7}{|c|}{\textbf{CHẤM BÀI VÀ SỬA BÀI KIỂM TRA CUỐI HỌC KỲ 2}} \\ \hline

% TUẦN 15 - 17
\textbf{15 - 17} & \multicolumn{6}{l|}{\textbf{Ôn tập cuối năm, ôn thi tốt nghiệp THPT.}} \\ \hline
\end{longtable}

\vspace{1.5cm}
% --- CHỮ KÝ ---
\noindent
\begin{minipage}[t]{0.45\textwidth}
    \centering
    \textbf{P. HIỆU TRƯỞNG}\\
    \vspace{2.5cm}
    \textbf{...........................................}
\end{minipage}
\hfill
\begin{minipage}[t]{0.45\textwidth}
    \centering
    \textit{Tp. HCM, ngày 05 tháng 09 năm 2025}\\
    \textbf{TỔ TRƯỞNG}\\
    \vspace{2cm}
    \textbf{Đỗ Thị Bạch Lan}
\end{minipage}

\end{document}
